% Companion section to cm_norms.tex; no document preamble.
% Requires the definitions and labels introduced in that section.

\section{Genus characters prove the recorded CM ratios}
\label{cmgenus:sec:main}

A class product and a class ratio contain different information.
The preceding norm formula controls the former. A nontrivial genus
character supplies the missing signed product and proves all three
recorded weight-four ratios at discriminants $-15,-20,-39$.
The same calculation gives a theorem for every even weight and three
explicit weight-six ratios.

The analytic mechanism is classical: twist the Kronecker limit formula
by a genus character and use Kronecker's factorization of its Hecke
$L$-function. A primary reference for the factorization, including its
precise character convention, is Murty and Murty,
Section~5~\cite{cmgenus:MurtyMurty}; see also Masri\cite[eqs.~(1.3)--(1.5)]{cmgenus:Masri}. We give the normalization argument
because confusing the ordinary and narrow real-quadratic class numbers,
or overlooking the units at discriminants $-3$ and $-4$, changes the
answer.

\subsection{A genus ratio theorem in arbitrary even weight}

Suppose that the fundamental discriminant $D<-4$ has a factorization
\[
 D=eD_-,\qquad e>1,\qquad D_-<0,
\]
where $e,D_-$ are coprime fundamental discriminants. Let
$\varepsilon:\mathrm{Cl}(D)\longrightarrow\{1,-1\}$ be the
corresponding nontrivial genus character. Concretely, if a form $Q$
represents an integer $n$ coprime to $D$, then
$\varepsilon(Q)=\chi_e(n)$. This value is independent of the represented
integer chosen. Retain the reduced representatives $\tau_Q$ and the
normalizations of Section~\ref{cmnorm:sec:main}, and put
\begin{align}
 A_\varepsilon&=\prod_{Q\in\mathcal Q_D}a_Q^{\varepsilon(Q)},&
 H_D^\pm(X)&=\prod_{\varepsilon(Q)=\pm1}(X-j(\tau_Q)),\label{cmgenus:eq:data}\\
 \lambda_-&=L(0,\chi_{D_-})
      =\frac{2h(D_-)}{\omega(D_-)},&
 \omega(D_-)&=|\mathcal O_{D_-}^{\times}|.\notag
\end{align}
Here $h(D_-)$ is the imaginary-quadratic class number. Let $h(e)$ be
the \emph{ordinary} real-quadratic class number, and let $u_e>1$ be the
least unit greater than one in the real embedding of
$\mathbb Q(\sqrt e)$. Its norm is allowed to be $-1$.
In particular,
\[
 \lambda_{-3}=\frac13,\qquad \lambda_{-4}=\frac12;
\]
replacing either of these values by its imaginary class number would
give an incorrect result.

Both genus factors are monic of degree $h(D)/2$, with
$H_D=H_D^+H_D^-$. The two sign sets are unions of genera; in the
three examples below they are the two genera themselves.
Their coefficients lie in the ring of integers of
$\mathbb Q(\sqrt e)$. Indeed, the kernel of $\varepsilon$ fixes the
quadratic genus extension $\mathbb Q(\sqrt D,\sqrt e)$ of
$\mathbb Q(\sqrt D)$; the two products are invariant under this kernel.
Complex conjugation inverts ideal classes, and inversion preserves
$\varepsilon$, so both products are real. Their coefficients therefore
belong to the real quadratic subfield. The coefficients are algebraic
integers because the singular moduli are algebraic integers.

\begin{theorem}[Genus ratios at every even weight]
\label{cmgenus:thm:allweights}
For even $w\geq4$, suppose that $G_w(\tau_Q)\neq0$ whenever
$\varepsilon(Q)=-1$. Then
\begin{equation}
 \prod_{Q\in\mathcal Q_D}|G_w(\tau_Q)|^{\varepsilon(Q)}
 =A_\varepsilon^{w/2}\,
   u_e^{-w h(e)\lambda_-}
   \left|
    \frac{\operatorname{Res}(H_D^+,F_w)}
         {\operatorname{Res}(H_D^-,F_w)}
   \right|^{1/12}.
 \label{cmgenus:eq:allweights}
\end{equation}
The polynomials $F_w$ and the resultant convention are those of
Lemma~\ref{cmnorm:lem:polynomial}. All real roots on the right are
positive roots; a vanishing numerator is permitted.
\end{theorem}

\begin{proof}
Write $d=|D|$ and
\[
 Z(\tau,s)=\sum_{(m,n)\neq(0,0)}
          \frac{(\operatorname{Im}\tau)^s}{|m+n\tau|^{2s}}.
\]
The quadratic-form and ideal-class zeta correspondence, with two units
because $D<-4$, gives
\begin{equation}
 \sum_Q\varepsilon(Q)Z(\tau_Q,s)
 =2\left(\frac{\sqrt d}{2}\right)^s
       L(s,\chi_e)L(s,\chi_{D_-}).
 \label{cmgenus:eq:twistedepstein}
\end{equation}
Here we used Kronecker's factorization
$L_{\mathbb Q(\sqrt D)}(s,\varepsilon)
=L(s,\chi_e)L(s,\chi_{D_-})$.
For completeness, at an unramified split prime the two quadratic
characters have the same sign and give the two equal prime-ideal
factors; at an inert prime their signs are opposite and their product
is $(1-p^{-2s})^{-1}$, the factor for the principal ideal $(p)$.
At a ramified prime, coprimality of $e,D_-$ leaves precisely the factor
of the unramified Dirichlet character. Thus the Euler products agree
in $\operatorname{Re}s>1$, and continuation gives the identity.

Since $\chi_e$ is nontrivial and even, $L(0,\chi_e)=0$.
Differentiate \eqref{cmgenus:eq:twistedepstein} at zero and use
\[
 Z'(\tau,0)=-2\log\bigl(2\pi\sqrt y\,|\eta(\tau)|^2\bigr),
 \qquad y=\operatorname{Im}\tau.
\]
The constant disappears because $\sum_Q\varepsilon(Q)=0$. We obtain
\begin{equation}
 \sum_Q\varepsilon(Q)\log\bigl(\sqrt{y_Q}
                                  |\eta(\tau_Q)|^2\bigr)
 =-\lambda_-L'(0,\chi_e).
 \label{cmgenus:eq:etaweighted}
\end{equation}
Since $y_Q=\sqrt d/(2a_Q)$, its signed logarithmic sum is
$\sum_Q\varepsilon(Q)\log y_Q=-\log A_\varepsilon$.
Consequently
\begin{equation}
 \prod_Q|\Delta(\tau_Q)|^{\varepsilon(Q)}
 =A_\varepsilon^6
       \exp\bigl(-12\lambda_-L'(0,\chi_e)\bigr).
 \label{cmgenus:eq:delta}
\end{equation}

To fix the real-quadratic normalization, the completed function
\[
 \Lambda_e(s)=(e/\pi)^{s/2}\Gamma(s/2)L(s,\chi_e)
\]
satisfies $\Lambda_e(s)=\Lambda_e(1-s)$.
Its values at zero and one give
$2L'(0,\chi_e)=\sqrt e\,L(1,\chi_e)$.
The real-quadratic class number formula, in the ordinary-class-number
and least-unit convention specified above, is
\[
 L(1,\chi_e)=\frac{2h(e)\log u_e}{\sqrt e}.
\]
Thus $L'(0,\chi_e)=h(e)\log u_e$, and
\eqref{cmgenus:eq:delta} becomes
\begin{equation}
 \boxed{\displaystyle
 \prod_Q|\Delta(\tau_Q)|^{\varepsilon(Q)}
 =A_\varepsilon^6 u_e^{-12h(e)\lambda_-}.}
 \label{cmgenus:eq:deltaunit}
\end{equation}
Finally, apply
$|E_w(\tau_Q)|^{12}
=|\Delta(\tau_Q)|^w|F_w(j(\tau_Q))|$ to the signed product.
The factors $2\zeta(w)$ cancel because the plus and minus character
fibers have the same cardinality. The remaining products of $F_w(j)$ are exactly the stated
resultants. Taking the positive twelfth root proves the formula.
\end{proof}

In particular, write
$R_w(D)=\prod_Q|G_w(\tau_Q)|^{\varepsilon(Q)}$. The two smallest
weights give the especially convenient identities
\begin{align}
 R_4(D)^3
 &=A_\varepsilon^6u_e^{-12h(e)\lambda_-}
       \left|\frac{H_D^+(0)}{H_D^-(0)}\right|,
       \label{cmgenus:eq:four}\\
 R_6(D)^2
 &=A_\varepsilon^6u_e^{-12h(e)\lambda_-}
       \left|\frac{H_D^+(1728)}{H_D^-(1728)}\right|.
       \label{cmgenus:eq:six}
\end{align}
Genus theory identifies the field of the factors $H_D^\pm$; these
equations additionally account for the periods of the raw modular
forms. Claiming a quadratic raw ratio requires verifying the root
extraction. We now perform that step exactly.

\subsection{The three recorded ratios, now theorems}

The required character and unit data are as follows. A plus or minus
sign attached to a form denotes its genus-character value.
\begin{center}
\begin{tabular}{@{}cclccc@{}}
\toprule
$D$&$(e,D_-)$&plus forms&$A_\varepsilon$&$u_e$&$\lambda_-$\\
\midrule
$-15$&$(5,-3)$&$(1,1,4)$&$1/2$&$(1+\sqrt5)/2$&$1/3$\\
$-20$&$(5,-4)$&$(1,0,5)$&$1/2$&$(1+\sqrt5)/2$&$1/2$\\
$-39$&$(13,-3)$&$(1,1,10),(3,3,4)$&$3/4$&$(3+\sqrt{13})/2$&$1/3$\\
\bottomrule
\end{tabular}
\end{center}
All the other reduced forms have minus sign. These signs can be
checked by represented integers: at $D=-15$, the nonprincipal form
represents $2$ and $\chi_5(2)=-1$; at $D=-20$, the nonprincipal form
represents $7$ and $\chi_5(7)=-1$; at $D=-39$, the forms with $a=2$
have $\chi_{13}(2)=-1$, whereas $(3,3,4)$ represents $10$ and
$\chi_{13}(10)=1$.
The class numbers $h(5)=h(13)=1$ follow from the real-quadratic
Minkowski bound $\sqrt e/2<2$.
The displayed units are fundamental: checking
$x^2-ey^2=\pm4$ below their sizes leaves no smaller positive unit.

The genus factors are explicitly
\begin{align}
 H_{-15}^{\pm}(X)
 &=X+\frac{191025\pm85995\sqrt5}{2},
       \label{cmgenus:eq:h15}\\
 H_{-20}^{\pm}(X)
 &=X-632000\mp282880\sqrt5,
       \label{cmgenus:eq:h20}\\
 H_{-39}^{\pm}(X)
 &=X^2+(165765798\pm45975573\sqrt{13})X\notag\\
 &\hspace{2em}+\frac{63399280527\pm17399806263\sqrt{13}}2.
       \label{cmgenus:eq:h39}
\end{align}
Multiplication gives the certified $H_D$ polynomials of
Section~\ref{cmnorm:sec:main}. The correspondence between these
quadratic-field factors and the displayed genera is also certified:
for each CM root, evaluate the other factor in a directed interval
enclosure and exclude zero. The attached receipt records the exact
rational interval endpoints, so the assignment is not a numerical
guess about which radical sign to choose.

\begin{theorem}[Exact weight-four genus ratios]
\label{cmgenus:thm:four}
With the character data in the table,
\begin{align}
 R_4(-15)&=\frac{(4+\sqrt5)^2}{44},
       \label{cmgenus:eq:r415}\\
 R_4(-20)&=\frac{29+12\sqrt5}{44},
       \label{cmgenus:eq:r420}\\
 R_4(-39)&=\frac{103299+4440\sqrt{13}}{181424}.
       \label{cmgenus:eq:r439}
\end{align}
The first equation is an absolute-value ratio; the precise real-part
version in the manuscript is
\begin{equation}
 \frac{G_4((1+i\sqrt{15})/2)}
      {\operatorname{Re}G_4((1+i\sqrt{15})/4)}
 =\frac{2(4+\sqrt5)^2}{77}.
 \label{cmgenus:eq:recorded15}
\end{equation}
For $D=-20$ and $D=-39$, the positive products represented by
$R_4(D)$ are exactly the products and ratios in the manuscript, without
an additional phase factor.
\end{theorem}

\begin{proof}
Substitute \eqref{cmgenus:eq:h15}--\eqref{cmgenus:eq:h39} and the table
into \eqref{cmgenus:eq:four}. In each of the two quadratic fields,
direct multiplication verifies that the cube of the claimed positive
number is the right-hand side. For example, for $D=-20$ this is the
explicit identity
\[
 \left(\frac{29+12\sqrt5}{44}\right)^3
 =\frac1{64}\left(\frac{1+\sqrt5}{2}\right)^{-6}
   \left|\frac{-632000-282880\sqrt5}
                  {-632000+282880\sqrt5}\right|.
\]
The identical finite verification with exponents $-4$ and $-4$, and
leading-coefficient ratios $1/2$ and $3/4$, proves the other two cases.
The attached exact-arithmetic script verifies all three equalities in
$\mathbb Q(\sqrt5)$ or $\mathbb Q(\sqrt{13})$ before any decimal
evaluation. Positivity makes the cube root unique.

For $D=-15$, the phase calculation of
Section~\ref{cmnorm:sec:phase} gives
$\operatorname{Re}G_4((1+i\sqrt{15})/4)
=\tfrac78|G_4((1+i\sqrt{15})/4)|$.
Multiplying \eqref{cmgenus:eq:r415} by $8/7$ proves
\eqref{cmgenus:eq:recorded15}; translation by one and reflection in the
imaginary axis account for the manuscript's positive-real-part
representatives. The phase discussion in the preceding section also
shows that the numerator and denominator genus products at $-20$ and
$-39$ are positive: boundary representatives contribute positive real
$G_4$, and the two interior representatives at $-39$ form a conjugate
pair.
\end{proof}

This supplies the missing proofs for the ratios originally displayed
in the manuscript's Results \texttt{cm:res:d15},
\texttt{cm:res:d20}, and \texttt{cm:res:genus}.
Combined with the proved total products, it also determines the
individual absolute values at class number two, and the absolute
products over each genus at class number four: if $T_w$ is the full
positive product, the two genus products are
$\sqrt{T_wR_w}$ and $\sqrt{T_w/R_w}$.

\subsection{Three explicit weight-six identities}

The same genus polynomials have further consequences beyond the
recorded weight-four data.

\begin{theorem}[Weight-six genus ratios]
\label{cmgenus:thm:six}
With the same genera and representatives,
\begin{align}
 R_6(-15)&=\frac{29+12\sqrt5}{88},
       \label{cmgenus:eq:r615}\\
 R_6(-20)&=\frac{601+252\sqrt5}{1672},
       \label{cmgenus:eq:r620}\\
 R_6(-39)&=\frac{1323+324\sqrt{13}}{1472}.
       \label{cmgenus:eq:r639}
\end{align}
The last two again hold for the corresponding raw genus products.
At discriminant $-15$ an explicit real-part identity is
\begin{equation}
 \frac{G_6((1+i\sqrt{15})/2)}
      {\operatorname{Re}G_6((1+i\sqrt{15})/4)}
 =\frac{2(29+12\sqrt5)}{121}.
 \label{cmgenus:eq:real6}
\end{equation}
\end{theorem}

\begin{proof}
Use \eqref{cmgenus:eq:six}. Exact arithmetic in the relevant quadratic
field shows that its right-hand sides are, respectively,
\begin{align*}
 &\frac{1561+696\sqrt5}{7744}
      =\left(\frac{29+12\sqrt5}{88}\right)^2,\\
 &\frac{678721+302904\sqrt5}{2795584}
      =\left(\frac{601+252\sqrt5}{1672}\right)^2,\\
 &\frac{3115017+857304\sqrt{13}}{2166784}
      =\left(\frac{1323+324\sqrt{13}}{1472}\right)^2.
\end{align*}
The roots on the right are positive. The valence formula gives the
only zero of $E_6$ in the standard fundamental domain at $i$; none of
the CM representatives under consideration is this point. Along the
boundary rays $\operatorname{Re}\tau=0$ and
$\operatorname{Re}\tau=-1/2$ above the relevant representatives,
$E_6$ is real, has no zero, and tends to one at infinity, so its values
are positive. Conjugate interior points again contribute a positive
product. This proves the raw-product assertion at $-20,-39$.

On the unit arc $|\tau|=1$ with
$\pi/2<\arg\tau\leq2\pi/3$, modularity and conjugation imply that
$\tau^3G_6(\tau)$ is real. It is positive at the hexagonal endpoint,
where $\tau^3=1$ and $G_6>0$, and cannot change sign before $i$.
Thus $G_6(\tau)=\tau^{-3}|G_6(\tau)|$ on this open arc. At
$\tau=(-1+i\sqrt{15})/4$,
\[
 \tau^{-3}=\frac{11+3i\sqrt{15}}{16},\qquad
 \operatorname{Re}G_6(\tau)=\frac{11}{16}|G_6(\tau)|.
\]
Reflection preserves the real part. Multiplying
\eqref{cmgenus:eq:r615} by $16/11$ proves
\eqref{cmgenus:eq:real6}.
\end{proof}

The quadratic norms provide useful exact checks:
\[
 N(R_4(-15))=N(R_4(-20))=\frac1{16},\qquad
 N(R_4(-39))=\frac{81}{256},
\]
and
\[
 N(R_6(-15))=N(R_6(-20))=\frac1{64},\qquad
 N(R_6(-39))=\frac{729}{4096}.
\]
These follow by conjugating $\sqrt e\mapsto-\sqrt e$ in the displayed
exact formulas. The identities are explicit consequences of the
classical genus-period mechanism; no assertion of historical priority
for individual numerical evaluations is intended.

% Insert after the weight-six genus-ratio subsection of cm_genus.tex.
% Uses the reduced modular polynomial of newmodular:eq:reduced.

\subsection{A fixed field contains all even-weight genus ratios}

The low-weight identities imply a uniform field-degree theorem.
For a fixed genus character, increasing the weight introduces no
additional algebraic field extension beyond the two seed ratios.

\begin{theorem}[Uniform algebraic degree of genus ratios]
\label{cmgenus:thm:degree}
Under the assumptions of Theorem~\ref{cmgenus:thm:allweights}, put
$K=\mathbb Q(\sqrt e)$ and
\[
 \mathcal K_{D,\varepsilon}=K(R_4(D),R_6(D)).
\]
Then
\[
 [\mathcal K_{D,\varepsilon}:K]\leq6,
 \qquad [\mathcal K_{D,\varepsilon}:\mathbb Q]\leq12,
\]
and every defined ratio $R_w(D)$ with even $w\geq4$ belongs to this
same field. The degree bound is independent of $w$ and of the imaginary
class number $h(D)$.
\end{theorem}

\begin{proof}
The genus period ratio
\[
 B_{D,\varepsilon}=A_\varepsilon^6
              u_e^{-12h(e)\lambda_-}
\]
belongs to $K$. Indeed,
\[
 12h(e)\lambda_-=
 \frac{24h(e)h(D_-)}{\omega(D_-)}\in\mathbb Z,
\]
because $\omega(D_-)$ is $2$, $4$, or $6$.
Equations~\eqref{cmgenus:eq:four} and \eqref{cmgenus:eq:six} therefore
give $R_4(D)^3\in K$ and $R_6(D)^2\in K$.
The seed ratios are nonzero: $E_4$ vanishes only at the elliptic orbit
of discriminant $-3$, and $E_6$ only at that of discriminant $-4$,
whereas $D<-4$ is fundamental. Adjoining the two positive seed ratios
has degree at most $3\cdot2=6$ over $K$.

To prove that this one field suffices, use the reduced modular
polynomial of \eqref{newmodular:eq:reduced}:
\begin{equation}
 E_w=E_4^aE_6^b\Delta^m P_w(j),\qquad
 a\in\{0,1,2\},\quad b\in\{0,1\},\quad
 w=4a+6b+12m,
 \label{cmgenus:eq:smallpolynomial}
\end{equation}
where $m\geq0$ and $P_w\in\mathbb Q[X]$ is monic of degree $m$.
The algebraic reason is short: the weight condition forces every
monomial of $E_w\in\mathbb Q[E_4,E_6]$ to have exponents
$a+3r,b+2s$ with $r+s=m$. Substituting
$E_4^3=\Delta j$ and $E_6^2=\Delta(j-1728)$ gives the displayed
form; the constant term $E_w=1+O(q)$ makes $P_w$ monic.

The signed absolute product gives the formula without fractional
root extraction
\begin{equation}
 R_w(D)=R_4(D)^aR_6(D)^b B_{D,\varepsilon}^{m}
       \left|\frac{\operatorname{Res}(H_D^+,P_w)}
                    {\operatorname{Res}(H_D^-,P_w)}\right|.
 \label{cmgenus:eq:noroot}
\end{equation}
The resultant ratio is in $K$, because both genus factors have
coefficients there. Its absolute value is also in $K$, which is a
real field. Since the seed forms and $\Delta$ never vanish at these
CM points, the denominator condition on $R_w$ guarantees a nonzero
denominator resultant. This proves the field assertion.
\end{proof}

\begin{corollary}[Quadratic closure for the recorded discriminants]
\label{cmgenus:cor:closure}
For $D=-15,-20,-39$, every defined genus ratio $R_w(D)$ with even
$w\geq4$ lies in $\mathbb Q(\sqrt5)$, $\mathbb Q(\sqrt5)$, or
$\mathbb Q(\sqrt{13})$, respectively. Each is computable in that
quadratic field using the reduced polynomial $P_w$, whose degree
is at most $\lfloor w/12\rfloor$.
\end{corollary}
\begin{proof}
Theorems~\ref{cmgenus:thm:four} and \ref{cmgenus:thm:six} place both
seed ratios in the indicated quadratic field. Formula
\eqref{cmgenus:eq:noroot} completes the proof.
\end{proof}

For example, $E_{14}=E_4^2E_6$ gives
$R_{14}(D)=R_4(D)^2R_6(D)$.
At weight $16$, $a=m=1,b=0$ and
$P_{16}(X)=X-3456000/3617$, so
\[
 R_{16}(D)=R_4(D)B_{D,\varepsilon}
    \left|\frac{H_D^+(3456000/3617)}
                 {H_D^-(3456000/3617)}\right|.
\]
The equal degrees of the two genus factors cancel the signs in the
linear-polynomial resultants. This formula, together with the three
explicit pairs $H_D^\pm$, generates further exact quadratic identities
using rational arithmetic alone.


\subsection{Certificates and further directions}

The companion \path{cm_genus.py} imports the directed-interval
Lambert-series evaluator from \mbox{\texttt{cm\_norms.py}}. It first certifies
the integral Hilbert polynomials, then verifies the displayed genus
factorizations by exact symbolic multiplication. For every reduced
form it records a represented integer witnessing the genus character,
and proves that the opposite genus factor excludes zero at its CM
$j$-value. It next checks the cubic and square identities exactly.
The receipt \path{cm_genus_receipt.json} separates these proof steps
from six ordinary 120-digit Fourier-series regressions. A second checker, \path{verify_genus_arithmetic.py}, independently verifies
all six radical equalities and their positive-root choices using only exact
standard-library rational arithmetic in the quadratic fields. The numerical regressions
have relative discrepancies below $3.1\times10^{-120}$ and are
independent consistency checks, not the justification for replacing
the manuscript's numerical evidence by theorems.

The genus theorem answers the three recorded ratio questions while
leaving several precise extensions open for this project.
\begin{enumerate}
\item Extend the genus-ratio certificate to nonmaximal quadratic
orders, including the local conductor factors in the twisted
Kronecker factorization. The untwisted conductor correction proved
earlier cannot simply be copied into a signed genus product.
\item Beyond the three discriminants of Corollary~\ref{cmgenus:cor:closure},
classify the discriminants and weights for which the raw ratio
$R_w(D)$ itself belongs to $\mathbb Q(\sqrt e)$, rather than merely
having a controlled algebraic power in that field. Exact root
extraction and representative phases must be part of the statement.
\item For class groups with several independent genus characters,
combine their signed formulas by finite Fourier inversion to recover
the product over every genus. This determines genus products but does
not distinguish classes inside one genus.
\item Investigate twists by nonreal class characters. Their Hecke
$L$-functions no longer reduce to a pair of quadratic Dirichlet
$L$-functions; the resulting period data measure exactly the
information absent from genus theory.
\end{enumerate}
